\section{Results}\label{sec:results}

\subsection{Test performance}

The results divide by model tier. Top-of-line models score \FlagScoreMin--\FlagScoreMax{} across all
conditions and both methods: for these models the method makes little practical difference. Mid-range models
vary more (\SmallScoreMin--\SmallScoreMax); here the state machine wins more often. \Cref{tab:matrix} gives the score of every run and \cref{fig:scores}
plots them; \cref{tab:pairs} compares the two methods pair by pair with 95\% confidence intervals. Cost and time are
in \cref{sec:cost}.

\begin{table}[tbp]\centering\small
\caption{Test AU-PRC (average precision) by model, method and reasoning effort. Each model's best run is marked in
bold; the best overall is marked bold and underlined. $^{\dagger}$: a repeated run (\cref{sec:repeats}).
$^{\ddagger}$: test predictions produced by hand from the run's selected model (\cref{sec:limitations}).}
\label{tab:matrix}
\MatrixTable
\end{table}

\begin{figure}[tbp]\centering
% Every run's test AU-PRC, one row per model (series from generated.tex).
\begin{tikzpicture}
\begin{axis}[width=0.82\linewidth, height=11.5cm, xmin=0.60, xmax=0.94, ymin=0.5, ymax=9.5,
  dotrows, yticklabel style={font=\small}, xlabel={test AU-PRC},
  xmajorgrids, grid style={black!8}, axis line style={black!40},
  xlabel style={yshift=-2pt},
  legend style={font=\footnotesize, at={(0.5,1.04)}, anchor=south, draw=black!20, fill=white, legend columns=4},
  legend cell align=left,
  every legend image/.append style={mark size=2.5pt},
  every axis plot/.append style={only marks, mark size=2.8pt}]
\DotPlots
\legend{{skill (\texttt{none})},{skill (\texttt{high})},{SM (\texttt{none})},{SM (\texttt{high})}}
\end{axis}
\end{tikzpicture}

\caption{Test AU-PRC of every run. Circles: skill; squares: state machine (Agentomics-ML). Filled: reasoning \texttt{high}; open: reasoning \texttt{none}. The bar through each marker is the run's 95\% bootstrap
interval (test-set sampling only).}
\label{fig:scores}
\end{figure}


\begin{table}[tbp]\centering\small
\caption{Paired comparison on the test set. SM: state machine (Agentomics-ML). $\Delta$: SM minus skill. 95\% CI:
percentile interval over 2,000 paired bootstrap resamples; $\Delta$ in bold where the interval excludes zero.
The interval reflects test-set sampling only, not run-to-run variation (\cref{sec:limitations}). $^{\dagger}$,
$^{\ddagger}$: as in \cref{tab:matrix}.}
\label{tab:pairs}
\PairTable\\[0.4em]
\PairTableNote
\end{table}

\paragraph{Score uncertainty.} Each run is evaluated on the same 1,018 held-out test examples, and those specific examples introduce sampling variation: a model that scores 0.80 on this test set might score differently on a different draw from the same distribution. To quantify this, we resampled the test set 2,000 times with replacement and recomputed AU-PRC each time; the central 95\% of that distribution gives a confidence interval for the run's score. For a typical run, this interval extends roughly \CIHalfMed{} in each direction, so a difference of that size or smaller between two rows of \cref{tab:matrix} may reflect the particular test examples rather than a genuine performance gap.

The paired comparison in \cref{tab:pairs} is more reliable for judging the two methods head to head: because both methods are scored on the same test examples, most of that sampling noise cancels in the difference. Across the \WilcoxonN{} model-and-reasoning pairs, the state machine is clearly better in \PairsSm{}, the skill in \PairsSkill{}, and \PairsTie{} are indistinguishable. Taken together, neither method has a clear overall edge (Wilcoxon signed-rank test, $p=\WilcoxonP$, median difference \MedianDiff{}): the state machine's wins on three mid-range models are offset by ties on top-of-line models and by the pairs the skill wins. With the original runs in place of the nine repeats the counts are \OrigPairsSm{}, \OrigPairsSkill{} and \OrigPairsTie{} ($p=\OrigWilcoxonP$; \cref{sec:repeats}).

\paragraph{Top-of-line models perform alike under both methods.} With reasoning
\texttt{high}, all \FlagHighTies{} top-of-line pairs are ties, the largest difference being \FlagHighMax. With reasoning
\texttt{none}, the state machine is ahead for Fable and Astra; the
skill catches up once the model reasons.

\paragraph{Mid-range models gain from the state machine, with exceptions.} Among the other five models the state
machine is clearly better in \SmallSm{} of \SmallPairs{} pairs (Haiku and DeepSeek at both settings, and Sonnet with reasoning \texttt{high}) and the
skill in \SmallSkill{}: Qwen 3.5 122B and GPT-4.1, both with reasoning \texttt{none}. The
remaining three pairs, Sonnet with reasoning \texttt{none} and Qwen and GPT-4.1 with reasoning \texttt{high}, are ties
(\cref{tab:pairs}; GPT-4.1 is discussed in \cref{sec:agentomics}).

Higher reasoning effort improved test AU-PRC in \HighWins{} of the \HighPairs{} model--method pairs, which is not distinguishable
from chance. The best single run was \BestCell, at \BestScore.

\subsection{Iterations}\label{sec:iterations}

\Cref{fig:itermeans} shows the mean test AU-PRC of the model a run would have submitted after one, two and three
iterations, for each tier, method and reasoning setting. \Cref{fig:iterations} shows every model separately, and \cref{tab:iterations} in \cref{app:tables} gives the
numbers, the iteration each run submitted and the iteration whose model scored best on test.

\begin{figure}[tbp]\centering
% Mean test AU-PRC after 1, 2 and 3 iterations, by tier, method and reasoning effort (series from generated.tex).
\begin{tikzpicture}
\begin{groupplot}[group style={group size=2 by 1, horizontal sep=1.6cm},
  width=6.6cm, height=5.6cm, xmin=0.7, xmax=3.3, xtick={1,2,3}, ymin=0.55, ymax=0.94,
  ytick={0.55,0.60,0.65,0.70,0.75,0.80,0.85,0.90},
  ytick scale label code/.code={}, scaled y ticks=false,
  yticklabel style={/pgf/number format/fixed, /pgf/number format/fixed zerofill, /pgf/number format/precision=2},
  title style={font=\small\bfseries}, axis line style={black!40}, label style={font=\footnotesize},
  ymajorgrids, grid style={black!8},
  every axis plot/.append style={mark size=2.6pt, line width=1.1pt}]
\IterMeanPlots
\end{groupplot}
\node[anchor=south] at ($(group c1r1.north)!0.5!(group c2r1.north)+(0,0.7cm)$) {\pgfplotslegendfromname{itermeanlegend}};
\end{tikzpicture}

\caption{Mean test AU-PRC of the model a run would have submitted after 1, 2 and 3 iterations, by model tier. Circles:
skill; squares: state machine (Agentomics-ML). Solid lines and filled markers: reasoning \texttt{high}; dashed lines
and open markers: reasoning \texttt{none}. Both panels share one y-axis; the dotted line is the human state of the art (0.860). A run with no working model yet is left out
of that point's mean, and a run with fewer iterations keeps what it ended with. \Cref{fig:iterations} shows every model separately.}
\label{fig:itermeans}
\end{figure}

\begin{figure}[tbp]\centering
% Test AU-PRC of the model each run would have submitted after 1, 2 and 3 iterations (series from generated.tex).
\begin{tikzpicture}
\begin{groupplot}[group style={group size=3 by 3, horizontal sep=1.2cm, vertical sep=1.45cm},
  width=5.0cm, height=4.3cm, xmin=0.7, xmax=3.3, xtick={1,2,3},
  enlarge y limits={abs=0.022},
  ytick scale label code/.code={}, scaled y ticks=false,
  yticklabel style={/pgf/number format/fixed, /pgf/number format/fixed zerofill, /pgf/number format/precision=2},
  title style={font=\footnotesize\bfseries, yshift=-3pt, text height=1.6ex, text depth=0.4ex}, axis line style={black!40},
  label style={font=\footnotesize}, xlabel style={yshift=2pt}, ylabel style={yshift=-3pt},
  ymajorgrids, grid style={black!8},
  every axis plot/.append style={mark size=2.1pt, line width=0.9pt}]
\IterPlots
\end{groupplot}
\node[anchor=south] at ($(group c2r1.north)+(0,0.75cm)$) {\pgfplotslegendfromname{iterlegend}};
\end{tikzpicture}

\caption{Test AU-PRC of the model each run would have submitted after 1, 2 and 3 iterations, one panel per model.
Circles: skill; squares: state machine (Agentomics-ML). Solid lines and filled markers: reasoning \texttt{high};
dashed lines and open markers: reasoning \texttt{none}. Dotted line: the human state of the art (0.860), where a
panel's range reaches it. Each panel has its own y-axis, so slopes compare within a panel, not across panels;
points are nudged sideways so that equal scores stay apart.}
\label{fig:iterations}
\end{figure}

Most of the gain arrives by the second iteration. Top-of-line models mostly start at 0.85--0.89 and add 0.01--0.04 afterwards; the exceptions are Astra and Gemini with
reasoning \texttt{none}, which start between 0.69 and 0.83, and Astra under the state machine with reasoning \texttt{high} (0.76). Mid-range models make their largest step at iteration 2
(Haiku rises from 0.44--0.58 to 0.67--0.74). Later iterations
can be by far the most expensive: DeepSeek under the state machine with reasoning \texttt{none} had 0.881 after seven
minutes, and its second and third iterations took a further 140 minutes without improving on it.

\paragraph{Single shot.} \Cref{tab:single} compares the two methods after the first iteration only, as if each run had
been given one attempt, pair by pair and with the same paired bootstrap as \cref{tab:pairs}. Seventeen of the eighteen pairs can be compared this way; the eighteenth is Gemini with reasoning \texttt{none},
whose skill run had no working model after its first iteration. Among the seventeen the state machine is the stronger
single shot: it is clearly ahead in \FirstCiSm{}, the skill in \FirstCiSkill{}, and \FirstCiTie{} are ties (Wilcoxon
signed-rank test, $p\FirstWilcoxonP$, median difference \FirstMedianDiff). The advantage is clearest among mid-range models:
\FirstMidCiSm{} of the \FirstMidCiPairs{} mid-range pairs, against \FirstFlagCiSm{} of the \FirstFlagCiPairs{}
top-of-line pairs. On average the state machine's first model scores \FirstSmMean{} and the skill's \FirstSkillMean{}.

The skill catches up once it is allowed to iterate. After one iteration the state machine leads by \FirstGap{} on
average; by the end its lead is \FinalGap. Mid-range models have the most to gain from iterating: their later
iterations add \FirstSmallGain{} on average, against \FirstFlagGain{} for top-of-line models. For many runs a single
attempt is already enough to be competitive: \FirstAboveSota{} of the \FirstRuns{} first models beat the human state
of the art.

\begin{table}[tbp]\centering\small
\caption{Single-shot comparison: test AU-PRC of the model each run had after its first iteration. SM: state machine
(Agentomics-ML). $\Delta$: SM minus skill, with a 95\% paired bootstrap interval as in \cref{tab:pairs}; bold where the
interval excludes zero. N/A: no working model after the first iteration. $^{\dagger}$, $^{\ddagger}$: as in
\cref{tab:matrix}.}
\label{tab:single}
\FirstPairTable
\end{table}

\paragraph{How many iterations.} \Cref{tab:budget} shows what each method has after one, two and three iterations,
and what share of its total cost and training time it has spent by then. Two iterations capture most of the result:
stopping there lowers the mean score by \TwoLoss{} and uses \TwoCostShare\% of the cost and \TwoTimeShare\% of the
training time. \TwoWithin{} of the \TwoRuns{} runs with a model after two iterations are already within 0.005 of their
final score.

\begin{table}[tbp]\centering\small
\caption{Iteration budget: the mean test AU-PRC each method has after one, two and three iterations (the model a run
would have submitted at that point), with the share of the method's total cost and training time spent by then. A run
with fewer iterations keeps what it ended with.}
\label{tab:budget}
\BudgetTable
\end{table}

\paragraph{A run does not always submit its best model.} Each run submits the iteration that scored highest on its
own validation data, since it never sees the test set. That is usually, but not always, the iteration that scores
highest on test: the two agree in \PickedBest{} of the \ScoredRuns{} runs. The clearest miss is GPT-4.1 under the state
machine with reasoning \texttt{high}. Its second iteration scored 0.752 on test, but its third scored slightly higher
on validation (0.721 against 0.705) and was submitted instead, at 0.726.
In the other runs where the two disagree, the cost is at most 0.020.

\subsection{Comparison with Agentomics-ML}\label{sec:agentomics}

\Cref{tab:agentomics} sets our GPT-4.1 runs beside those published for Agentomics-ML \citep{martinek2025agentomics} on
the same dataset and the same train/test split, with average precision, the metric used here. We make this comparison
for GPT-4.1 only, the model the published results use: the other models here are newer, and setting their scores
against the published ones would measure the models, not the method.

\begin{table}[tbp]\centering\small
\caption{Test AU-PRC (average precision) of GPT-4.1 on AGO2 CLASH. Agentomics-ML as reported by
\citet{martinek2025agentomics}, arXiv v1, over five runs per configuration. This study: two runs per method.}
\label{tab:agentomics}
\begin{tabular}{@{}l l r r@{}}\toprule
\textbf{Method} & \textbf{Model} & \textbf{Mean} & \textbf{Best}\\\midrule
Agentomics-ML, no feedback & GPT-4.1 & 0.698 & 0.724\\
Agentomics-ML & GPT-4.1 & 0.726 & 0.778\\\midrule
State machine (this study) & GPT-4.1, 2 runs & \GptSmMean & \GptSmBest\\
Skill (this study) & GPT-4.1, 2 runs & \GptSkillMean & \GptSkillBest\\
\bottomrule\end{tabular}
\end{table}

\paragraph{GPT-4.1 under both methods.} GPT-4.1 is the one model both studies share; all four of its runs here are
repeats made with Agentomics-ML 1.2.0 and the second skill agent (\cref{sec:repeats}). Under the state machine it scored
\GptSmNone{} and \GptSmHigh{} (two runs that are in effect repeats, \cref{sec:models}); in the first, the third iteration
failed its checks and the second was submitted. \GptSmHigh{} equals the mean Agentomics-ML reports for GPT-4.1 (means
of 0.698 and 0.726, best run 0.778) and \GptSmNone{} falls just below it, so our runs are broadly consistent with the
published result on the one model the two share.

Under the skill it scored \GptSkillNone{} and \GptSkillHigh{}, completing all three iterations in both runs for
\$5.92 and \$0.33, against \$1.50 and \$2.10 with the state machine. The skill is ahead in both pairs; the interval
of the pair at reasoning \texttt{high} includes zero (\cref{tab:pairs}). The oldest model here therefore did not need
the harness on this task: it followed the written procedure to the end and scored at least as well as with enforced
steps.

\subsection{Cost and time}\label{sec:cost}

The skill used \SkillCost{} and \SkillTime{} over its eighteen runs, \CostRatio{} of the cost and \TimeRatio{} of the
time of the state machine (\SmCost{} and \SmTime). The difference comes mostly from the number of model calls
(\SkillCalls{} against \SmCalls): the state machine starts a fresh model turn at every step, each with its own context,
whereas the agent works through the task in one turn. One state-machine record, Fable with reasoning \texttt{high},
sums two runs and accounts for \FableSmHighCost{} of that total (\cref{sec:limitations}). \Cref{fig:time} shows the cost and training time of every run and how the time divides across iterations;
\cref{tab:cost,tab:time} give the numbers.

\begin{table}[tbp]\centering\small
\caption{Cost per run in USD, as billed by OpenRouter. Each model's cheapest run is marked in bold; the cheapest overall
is marked bold and underlined.}
\label{tab:cost}
\CostTable
\end{table}

\begin{table}[tbp]\centering\small
\caption{Training time per run in minutes: model time plus code time up to the end of the last iteration (\cref{sec:accounting}). Each model's fastest run is marked in
bold; the fastest overall is marked bold and underlined.}
\label{tab:time}
\TimeTable
\end{table}

\begin{figure}[tbp]\centering
% Cost and training time per run on log axes, and each run's split of time across its iterations (series from generated.tex).
\begin{tikzpicture}
\pgfplotsset{runrows/.style={height=18.9cm, ymin=-0.7, ymax=35.7, axis line style={black!40},
  label style={font=\footnotesize}, tick label style={font=\scriptsize},
  extra y ticks={3.5,7.5,11.5,15.5,19.5,23.5,27.5,31.5}, extra y tick labels={,,,,,,,},
  extra y tick style={grid=major, grid style={draw=black!30, line width=0.5pt, dashed}, tick style={draw=none}}}}
\begin{axis}[name=cost, runrows, xmode=log, log basis x=10, width=0.33\linewidth, xmin=0.02, xmax=1500,
  xtick={0.1,1,10,100}, xticklabels={0.1,1,10,100}, timerows, yticklabel style={font=\scriptsize},
  xlabel={cost, USD (log)}, xmajorgrids, grid style={black!8}, every axis plot/.append style={mark size=2.2pt},
  legend to name=runlegend, legend columns=4, legend style={font=\scriptsize, draw=black!20, /tikz/every even column/.append style={column sep=0.5em}}]
\CostTotalPlots
\legend{{skill (\texttt{none})},{skill (\texttt{high})},{SM (\texttt{none})},{SM (\texttt{high})}}
\end{axis}
\begin{axis}[name=totals, at={(cost.east)}, anchor=west, xshift=0.35cm, runrows, xmode=log, log basis x=10, width=0.33\linewidth,
  xmin=0.5, xmax=900, xtick={1,10,100}, xticklabels={1,10,100}, ytick={0,...,35}, yticklabels={}, ytick style={draw=none},
  xlabel={time, min (log)}, xmajorgrids, grid style={black!8}, every axis plot/.append style={mark size=2.2pt}]
\TimeTotalPlots
\end{axis}
\begin{axis}[name=share, at={(totals.east)}, anchor=west, xshift=0.35cm, runrows, xbar stacked, bar width=7.5pt, width=0.29\linewidth,
  xmin=0, xmax=100, xtick={0,50,100}, ytick={0,...,35}, yticklabels={}, ytick style={draw=none},
  xlabel={time by iteration (\%)},
  legend style={font=\scriptsize, at={(0.5,1.008)}, anchor=south, legend columns=3, draw=black!20, /tikz/every even column/.append style={column sep=0.25em}}]
\TimeSharePlots
\legend{iteration 1, 2, 3}
\end{axis}
\node[anchor=south] at ($(cost.north)!0.5!(totals.north)+(-0.55cm,0.1cm)$) {\pgfplotslegendfromname{runlegend}};
\end{tikzpicture}

\caption{Cost and training time of every run. Runs group by model, four per model (skill or state machine, reasoning
\texttt{none} or \texttt{high}). Left: cost in USD; middle: training time in minutes; both on log scales, with the
value beside each marker (circles: skill, squares: state machine, filled: reasoning \texttt{high}). Right: how the
training time divides across iterations 1, 2 and 3. The state machine records the iteration it is in; for a skill run
an iteration is taken to end when the agent first writes into the next iteration's folder, so its split is
approximate. DeepSeek's state-machine run with reasoning \texttt{high} stopped during its second iteration, and Gemini's skill run with
reasoning \texttt{none} took five, the last three drawn together.}
\label{fig:time}
\end{figure}
